% Generated by scripts/build-user-guides.py. Edit the matching Markdown guide.
\documentclass[10pt,a4paper]{article}
\usepackage[margin=1.8cm,headheight=14pt]{geometry}
\usepackage[T1]{fontenc}
\usepackage{lmodern}
\usepackage{xcolor,hyperref,fancyhdr,enumitem,titlesec}
\definecolor{ExamPulseRed}{HTML}{E4002B}
\hypersetup{colorlinks=true,urlcolor=ExamPulseRed,linkcolor=ExamPulseRed}
\pagestyle{fancy}\fancyhf{}\lhead{\textbf{ExamPulse}}\rhead{Invigilator Guide}
\lfoot{\footnotesize 25 October 2026 | exampulse.xyz}\rfoot{\thepage}
\setlength{\parindent}{0pt}\setlength{\parskip}{5pt}
\titleformat{\section}{\large\bfseries\color{ExamPulseRed}}{}{0pt}{}
\titlespacing*{\section}{0pt}{8pt}{4pt}
\setlist[itemize]{leftmargin=*,nosep,after=\vspace{4pt}}
\begin{document}
{\LARGE\bfseries Invigilator Quick Guide}\par\medskip
Updated 25 October 2026. Sign in, choose your room, and mark attendance with ExamPulse.\par
\section*{1. Before the Exam}
Open \href{https://exampulse.xyz/scan}{exampulse.xyz/scan} on your phone before the exam starts. Have a charged phone, internet access, and your personal access code ready.\par
\textbf{New to ExamPulse?} Your exam administrator will provide your code separately. Look for an email titled \textbf{ExamPulse access code}, or a code shared with you privately. The assignment email tells you your exam and rooms; it does not include a new code. If you have not received a code, ask your administrator before the exam.\par
\textbf{Used it before?} Use your existing code unless your administrator has given you a replacement. If you cannot find it or it does not work, ask your administrator for help.\par
Keep your code private. Do not use someone else's code.\par
\section*{2. Sign In and Choose Your Room}
\begin{itemize}\item[1.] Open the link and enter your access code.\end{itemize}
\begin{itemize}\item[2.] Choose your assigned room. Check the exam name and room number against where you are working.\end{itemize}
\begin{itemize}\item[3.] Allow camera access when asked. Keep the page open while the scanner gets ready.\end{itemize}
First use can take around \textbf{20-40 seconds}, sometimes longer. You can choose \textbf{Use manual entry} and type a student number while you wait. You still need internet access.\par
\textbf{Room missing or incorrect?} Contact the coordinator. Do not choose another room just to continue.\par
\section*{3. Scan, Check, and Mark Present}
\begin{itemize}\item[1.] Hold the student's ID card steady. Put only the \textbf{printed student number} inside the red scan box, not the barcode.\end{itemize}
\begin{itemize}\item[2.] Check the student number, displayed student details, and room in the review.\end{itemize}
\begin{itemize}\item[3.] If correct, tap \textbf{Mark Present} and wait for the result. \textbf{Ready to mark} means you still need to tap the button; attendance has not been recorded yet.\end{itemize}
\begin{itemize}\item[4.] If the number is wrong, correct it and tap \textbf{Lookup Edited ID}. Check the new details before marking.\end{itemize}
If reading the card is difficult, use the torch control if available, or \textbf{Manual Mode}.\par
\section*{4. Type a Student Number Instead}
Choose \textbf{Manual Mode} or type into \textbf{Student number}, then tap \textbf{Lookup}. Check the details and tap \textbf{Mark Present}, just as you would after a camera scan.\par
If the student has no usable ID card, add the student number manually after checking the required form of Photo ID as specified by the coordinator. Add a brief note in \textbf{Comments}.\par
\newpage
\section*{5. If the Student Is in the Wrong Room}
Tell your room Team Leader and follow their instructions. The app shows the student's expected room and zone.\par
\begin{itemize}\item[\textbullet] Tap \textbf{Send to [expected room]} if the student should move. This does not mark them present in your room.\end{itemize}
\begin{itemize}\item[\textbullet] Tap \textbf{Mark present in [current room]} only if the Team Leader authorizes it.\end{itemize}
Check the room numbers on the buttons. Neither action changes the student's assigned room.\par
\section*{6. Other Results You May See}
\begin{itemize}\item[\textbullet] \textbf{Already marked:} do not mark again. Ask the Team Leader if the record needs checking.\end{itemize}
\begin{itemize}\item[\textbullet] \textbf{Student not found:} check the digits and look up the corrected number. If still not found, ask the Team Leader for help. Never use another student's number.\end{itemize}
\section*{7. If You Need Help}
\textbf{Scanner will not start:} use manual entry. If \textbf{Retry OCR} is available, tap it once after checking your connection. You do not need to wait for scanning to use manual entry.\par
\textbf{Camera will not open:} allow camera access in your browser settings. Tap \textbf{Restart Camera} if offered, or use manual entry.\par
\textbf{No internet, or the page closes:} tell the Team Leader. Use the approved paper process or a spare phone for new students, as directed. Keep the original phone for checking any uncertain entries. Typing a number still requires internet access.\par
\textbf{Code will not work:} check you entered your own code correctly. If it still fails, ask the coordinator for help.\par
\section*{8. Only If an Entry Is Waiting to Be Saved}
After a connection problem, you may see \textbf{Attendance sync queue}. If you do not see it, there is nothing extra to do.\par
\begin{itemize}\item[\textbullet] If it says \textbf{Pending synchronization} or \textbf{Synchronizing}, keep the page open with internet access and let it finish.\end{itemize}
\begin{itemize}\item[\textbullet] If it says \textbf{Could not synchronize} or \textbf{Conflict needs review}, show the coordinator. Do not repeatedly tap \textbf{Retry} or tap \textbf{Acknowledge} just to dismiss it.\end{itemize}
\begin{itemize}\item[\textbullet] Do not clear browser data or enter the same uncertain mark on another phone until the coordinator has checked it.\end{itemize}
\section*{9. Before You Leave}
Tell the Team Leader when attendance collection is finished. If any entries are still waiting to be saved, ask the coordinator to check them before you sign out or leave. Hand over any paper records.\par
\textbf{Need help?} Contact your room Team Leader or exam coordinator first. Application support: \href{mailto:ahmed.sohair.khan@rmit.edu.au}{ahmed.sohair.khan@rmit.edu.au}. Keep codes and student details out of group chats.\par
\end{document}
